\section{Tích lũy tư bản}

\subsection{Bản chất tích lũy tư bản}
Về bản chất, tích lũy tư bản là quá trình tư bản hóa giá trị thặng dư. Nhà tư bản không tiêu dùng toàn bộ giá trị thặng dư cho sinh hoạt cá nhân mà chuyển hóa một bộ phận của nó thành tư bản phụ thêm để mở rộng sản xuất. Do đó khi thị trường thuận lợi, nhà tư bản bán được hàng hóa, giá trị thặng dư sẽ ngày càng nhiều, nhà tư bản trở nên giàu có hơn.
\begin{itemize}
    \item \textbf{Tính chất hai mặt:} Một mặt, tích lũy là sự mở rộng quy mô sản xuất thực thể. Mặt khác, nó là quá trình vật hóa thời gian lao động không công của người công nhân thành phương tiện để bóc lột chính họ trong chu kỳ tiếp theo. Nguồn gốc duy nhất của tích lũy là giá trị thặng dư, và mục đích của nó là biến giá trị thặng dư thành tư bản.
    \item \textbf{Động lực khách quan:}
    \begin{itemize}
        \item Quy luật giá trị thặng dư: Sự khao khát giá trị thặng dư là vô tận, thúc đẩy nhà tư bản không ngừng gia tăng quy mô chiếm đoạt.
        \item Quy luật cạnh tranh: Trong guồng quay khốc liệt của thị trường, tích lũy là vũ khí sinh tồn. Nhà tư bản buộc phải cải tiến kỹ thuật và mở rộng quy mô để hạ giá thành, nếu không muốn bị đào thải hoặc thôn tính.
    \end{itemize}
\end{itemize}

\subsection{Nhân tố ảnh hưởng quy mô tích lũy}

Quy mô tích lũy không cố định mà biến thiên theo tỷ lệ phân chia giữa tích lũy và tiêu dùng. Với một tỷ lệ phân chia nhất định, quy mô này chịu sự chi phối của bốn nhân tố mang tính cấu trúc sau:

\begin{table}[H]
\centering
\renewcommand{\arraystretch}{1.5} % Tạo khoảng trống cho chữ dễ đọc
\setlength{\arrayrulewidth}{0.5pt}

\begin{tabularx}{\textwidth}{| >{\raggedright\arraybackslash\bfseries}p{3.5cm} | >{\raggedright\arraybackslash}X | >{\raggedright\arraybackslash}X |}
\hline
\rowcolor[HTML]{EFEFEF} 
Nhân tố quyết định & Cơ chế tác động kinh tế - chính trị & Hệ quả đối với quy mô tích lũy \\ \hline

Nâng cao mức độ khai thác sức lao động & 
Cắt giảm tiền lương thực tế xuống dưới giá trị sức lao động hoặc cưỡng bức kéo dài ngày lao động (tăng $m$ tuyệt đối). & 
Gia tăng trực tiếp khối lượng $m$ bằng cách chiếm đoạt trắng trợn quỹ sinh hoạt của người lao động. \\ \hline

Tăng năng suất lao động xã hội & 
Cải tiến kỹ thuật làm giảm giá trị tư liệu sinh hoạt, hạ thấp giá trị sức lao động ($v$), từ đó tăng tỷ suất bóc lột $m'$ (tăng $m$ tương đối). & 
Tăng quy mô tích lũy bằng cách cho phép nhà tư bản mua được nhiều tư liệu sản xuất và sức lao động hơn với cùng một lượng tiền. \\ \hline

Chênh lệch giữa tư bản sử dụng và tư bản tiêu dùng & 
Máy móc phục vụ với toàn bộ hiệu suất nhưng chỉ chuyển dần giá trị (khấu hao) vào sản phẩm. Phần giá trị phục vụ miễn phí này hoạt động như một lực lượng tự nhiên. & 
Tận dụng lực lượng sản xuất miễn phí từ máy móc để tăng quy mô tích lũy mà không cần tăng chi phí tư bản ứng trước tương ứng. \\ \hline

Đại lượng tư bản ứng trước & 
Với tỷ suất bóc lột $m'$ không đổi, khối lượng giá trị thặng dư $M$ tỷ lệ thuận với tổng số tư bản khả biến được đầu tư ban đầu. & 
Tư bản càng lớn, khối lượng $m$ thu được càng khổng lồ, tạo tiền đề cho những bước nhảy vọt về quy mô tích lũy tiếp theo. \\ \hline
\end{tabularx}
\caption{Các nhân tố quyết định quy mô tích lũy tư bản}
\end{table}

\subsection{Hệ quả của tích lũy tư bản}

\subsubsection{Tăng cấu tạo hữu cơ của tư bản}
“Cấu tạo hữu cơ của tư bản là cấu tạo giá trị được quyết định bởi cấu tạo kỹ thuật và phản ánh sự biến đổi của cấu tạo kỹ thuật của tư bản", (ký hiệu là c/v) \cite[trang 97]{gt}. Xu thế tất yếu của tích lũy là thay thế lao động sống bằng lao động vật hóa (máy móc) nhằm tăng năng suất và kiểm soát quá trình sản xuất. Khi cấu tạo hữu cơ $(c/v)$ tăng lên, tốc độ tăng của tư bản bất biến $(c)$ luôn nhanh hơn tư bản khả biến $(v)$. Hệ quả là nhu cầu tương đối về sức lao động giảm xuống, tạo ra "đội quân thất nghiệp dự bị" – một quần chúng lao động dư thừa thường xuyên phục vụ cho nhu cầu giãn nở của tư bản.

\subsubsection{Tăng tích tụ và tập trung tư bản}

Tích tụ tư bản là gia tăng quy mô tư bản cá biệt bằng cách tư bản hóa giá trị thặng dư tại doanh nghiệp, phản ánh mối quan hệ giữa giai cấp tư sản và công nhân. Tập trung tư bản là hợp nhất các tư bản có sẵn thông qua cạnh tranh và tín dụng. Tín dụng đóng vai trò là đòn bẩy mạnh mẽ nhất, cho phép tập hợp những nguồn vốn phân tán để triển khai các dự án quy mô khổng lồ, làm thay đổi nhanh chóng diện mạo cấu trúc kinh tế. Tích tụ và tập trung tư bản đều góp phần tạo tiền đề để có thể thu được nhiều giá trị thặng dư hơn cho người mua hàng hóa sức lao động \cite[trang 98]{gt}.

\subsubsection{Tăng chênh lệch giữa thu nhập của nhà tư bản và người lao động làm thuê cả tuyệt đối lẫn tương đối}

Tích lũy ở cực này (nhà tư bản) là tích lũy sự giàu sang, trong khi ở cực kia (người lao động) là sự tích lũy bần cùng, tạo nên hiện tượng bần cùng hóa người lao động, bao gồm:

\begin{itemize}
    \item Bần cùng hóa tuyệt đối: Sự sụt giảm mức sống thực tế trong các điều kiện khủng hoảng hoặc bóc lột cực đoan.
    \item Bần cùng hóa tương đối: Phần chia của giai cấp công nhân trong tổng giá trị mới tạo ra $(v+m)$ ngày càng nhỏ đi so với phần giá trị thặng dư $(m)$ mà nhà tư bản chiếm hữu, làm sâu sắc thêm hố ngăn cách xã hội.

\end{itemize}
